% Capítulo "Moradia" do relatório -- EDITÁVEL À MÃO (specs/2026-11-05_versao_final_site R5).
% Exportado uma vez por gera_latex.py --exportar-capitulos (2026-11-05); desde então o TEXTO é escrito aqui, não
% vem mais de relatorio/textos_curados.json (que é só do site). Os DADOS continuam gerados: as figuras
% (\figura{...}{../../visualizacoes/a4/...} e mapas/a4/) saem do analise.py; as tabelas no corpo
% (\input{gerado/tabelas_corpo/...}) e os rótulos do apêndice (\ref{tab:...}) saem do gera_latex.py, que confere
% a cada execução se tudo o que este arquivo cita existe.

\chapter{Moradia}\label{cap:eixo-7}

\begin{achados}\begin{itemize}

\item Metade das crianças até 72 meses do CadÚnico (51,6\%, ou 101.477) vive em domicílio com mais de 2 pessoas por dormitório

\item 43,2\% (86.393) vivem em domicílio em déficit habitacional, quase todo pelo peso do aluguel: 41,8\% estão em famílias de baixa renda que gastam mais de 30\% dela com aluguel

\item A inadequação da infraestrutura ou da construção atinge 5,9\% (11.589), com os maiores números em Santa Cruz (1.068) e Guaratiba (770)

\item A falta de banheiro (0,3\%) e de água canalizada (0,6\%) é rara

\end{itemize}\end{achados}

Os dados de moradia vêm do Cadastro Único e cobrem, portanto, as famílias de menor renda, não toda a cidade. A classificação segue a metodologia da Fundação João Pinheiro, adaptada às variáveis do cadastro: o déficit habitacional mede a necessidade de novas moradias (domicílios improvisados, rústicos, coabitação e ônus excessivo com aluguel), e a inadequação mede problemas no domicílio existente (infraestrutura e construção). O adensamento excessivo considera mais de 2 pessoas por dormitório.

\section{Crianças no CadÚnico em domicílios com inadequação habitacional}

\figura{grafico}{../../visualizacoes/a4/cadunico_criancas_inadequacao_componentes.pdf}{CADÚNICO: Crianças até 72 meses em domicílio com inadequação habitacional, por componente}{Elaboração IPP com dados de \citeonline{mds_cadunico}.}{graf:cadunico-criancas-inadequacao-componentes}

Segundo a aproximação da metodologia da Fundação João Pinheiro, 11.589 crianças até 72 meses do CadÚnico (5,9\%) vivem em domicílio com alguma inadequação habitacional. A maior parte se deve à infraestrutura (9.841 crianças, 5,0\%), sobretudo ao esgoto inadequado (4.676), à água sem canalização ou fora da rede geral (3.584) e ao lixo queimado, enterrado ou jogado (2.756). Os problemas da própria construção (inadequação edilícia) atingem 2.234 crianças (1,1\%): todos os cômodos usados como dormitório (1.018), piso de terra (859) e ausência de banheiro (627). Um domicílio pode ter mais de um componente.

\figura{mapa}{../../mapas/a4/mapa_cadunico_inadequacao_n_bairro_2026.pdf}{Crianças até 72 meses no CadÚnico em domicílio com inadequação habitacional, por bairro (número)}{Elaboração IPP com dados de \citeonline{mds_cadunico}. Limites de bairros: \citeonline{ipp_limites_bairros}. Sistema de referência SIRGAS 2000.}{mapa:mapa-cadunico-inadequacao-n-bairro-2026}

Em números absolutos, as crianças em domicílios com inadequação habitacional se concentram na Zona Oeste: Santa Cruz (1.068), Guaratiba (770) e Campo Grande (430) lideram, seguidos da Penha (308), na Zona Norte, e de Paciência (276). Os quatro primeiros são os únicos bairros com mais de 300 crianças nessa situação. Bairros com menos de 20 casos aparecem sem cor: estão somados por Região Administrativa ou Área de Planejamento na tabela.

\figura{mapa}{../../mapas/a4/mapa_percentual_cadunico_inadequacao_bairro_2026.pdf}{\% de crianças até 72 meses no CadÚnico em domicílio com inadequação habitacional, por bairro}{Elaboração IPP com dados de \citeonline{mds_cadunico}. Limites de bairros: \citeonline{ipp_limites_bairros}. Sistema de referência SIRGAS 2000. Nota: Escala de cor limitada a 20,8 (valores acima com a cor máxima); o valor real está na tabela do apêndice.}{mapa:mapa-percentual-cadunico-inadequacao-bairro-2026}

Em proporção, o caso mais destacado é o Alto da Boa Vista, onde 44,1\% das crianças do CadÚnico (83 de 188) vivem em domicílio com inadequação habitacional, mais que o dobro de qualquer outro bairro; por isso ele aparece na cor máxima e fora da escala. Seguem Vargem Grande (20,8\%), Barra de Guaratiba (17,7\%), Portuguesa (16,5\%) e Penha (15,1\%). Santa Cruz e Guaratiba, que têm os maiores números, ficam em 9,3\% e 11,6\%, acima da média do município (5,9\%).

\vertabelas{Tabela~\ref{tab:cadunico-inadequacao-componentes-0-a-5-2026}, Tabela~\ref{tab:tabela-mapa-cadunico-inadequacao-bairro-2026}; ver Apêndice~\ref{ap:eixo-7}}

\section{Crianças no CadÚnico em domicílios com adensamento habitacional excessivo (mais de 2 pessoas por dormitório)}

\figura{mapa}{../../mapas/a4/mapa_cadunico_adensamento_n_bairro_2026.pdf}{Crianças até 72 meses no CadÚnico em domicílio com adensamento excessivo, por bairro (número)}{Elaboração IPP com dados de \citeonline{mds_cadunico}. Limites de bairros: \citeonline{ipp_limites_bairros}. Sistema de referência SIRGAS 2000.}{mapa:mapa-cadunico-adensamento-n-bairro-2026}

O adensamento excessivo (mais de 2 pessoas por dormitório) é o problema de moradia mais comum entre as crianças do CadÚnico, e os maiores números estão na Zona Oeste: Santa Cruz (5.681 crianças), Campo Grande (5.249), Guaratiba (3.661), Bangu (3.053) e Realengo (2.877). Como atinge cerca de metade das crianças cadastradas, o mapa de número acompanha de perto a distribuição das crianças do CadÚnico pela cidade.

\figura{mapa}{../../mapas/a4/mapa_percentual_cadunico_adensamento_bairro_2026.pdf}{\% de crianças até 72 meses no CadÚnico em domicílio com adensamento excessivo, por bairro}{Elaboração IPP com dados de \citeonline{mds_cadunico}. Limites de bairros: \citeonline{ipp_limites_bairros}. Sistema de referência SIRGAS 2000. Nota: Escala de cor limitada ao percentil 95 (56,5); o valor real está na tabela do apêndice.}{mapa:mapa-percentual-cadunico-adensamento-bairro-2026}

Em 2026, 101.477 crianças até 72 meses do CadÚnico (51,6\%) vivem em domicílio com mais de 2 pessoas por dormitório. O percentual é alto em toda a cidade: mesmo os menores valores ficam acima de 30\%, como em Todos os Santos (32,3\%), Rocha (33,2\%) e Abolição (33,9\%). Os maiores estão em Vaz Lobo (61,6\%), Glória (57,9\%), São Conrado (57,8\%), Cacuia (57,5\%) e Sepetiba (56,8\%).

\vertabelas{Tabela~\ref{tab:tabela-mapa-cadunico-adensamento-bairro-2026}; ver Apêndice~\ref{ap:eixo-7}}

\section{Indicadores agregados de moradia (inadequação, saneamento, melhorias habitacionais)}

\figura{grafico}{../../visualizacoes/a4/cadunico_criancas_deficit_componentes.pdf}{CADÚNICO: Crianças até 72 meses em domicílio em déficit habitacional, por componente}{Elaboração IPP com dados de \citeonline{mds_cadunico}.}{graf:cadunico-criancas-deficit-componentes}

Pela aproximação da metodologia da Fundação João Pinheiro, 86.393 crianças até 72 meses do CadÚnico (43,2\%) vivem em domicílio em déficit habitacional. Quase todo o déficit vem do ônus excessivo com aluguel, quando famílias com renda de até 3 salários mínimos gastam mais de 30\% dela com aluguel: são 82.081 crianças (41,8\%). Os demais componentes são bem menores: domicílio improvisado (3.436 crianças), rústico (591) e coabitação (588). Um domicílio pode ter mais de um componente.

% Gerado por gera_latex.py -- não editar à mão.
\begin{landscape}
{\scriptsize\setlength{\tabcolsep}{3pt}
\begin{longtable}{@{}>{\raggedright\arraybackslash}p{3.6cm}rrrrrr@{}}
\caption{Crianças de 0 a 5 anos e famílias no CadÚnico, por indicador de moradia (Fundação João Pinheiro e adensamento)}\label{tab:cadunico-moradia-resumo-0-a-5-2026}\\
\toprule \textbf{Indicador} & \textbf{Crianças} & \parbox[b]{2.20cm}{\raggedleft\hyphenpenalty=10000\exhyphenpenalty=10000 \textbf{Base (crianças)}} & \parbox[b]{2.20cm}{\raggedleft\hyphenpenalty=10000\exhyphenpenalty=10000 \textbf{\% das crianças}} & \textbf{Famílias} & \parbox[b]{2.20cm}{\raggedleft\hyphenpenalty=10000\exhyphenpenalty=10000 \textbf{Base (famílias)}} & \parbox[b]{2.20cm}{\raggedleft\hyphenpenalty=10000\exhyphenpenalty=10000 \textbf{\% das famílias}} \\ \midrule \endfirsthead
\multicolumn{7}{@{}l}{\footnotesize\itshape (continuação)}\\ \toprule \textbf{Indicador} & \textbf{Crianças} & \parbox[b]{2.20cm}{\raggedleft\hyphenpenalty=10000\exhyphenpenalty=10000 \textbf{Base (crianças)}} & \parbox[b]{2.20cm}{\raggedleft\hyphenpenalty=10000\exhyphenpenalty=10000 \textbf{\% das crianças}} & \textbf{Famílias} & \parbox[b]{2.20cm}{\raggedleft\hyphenpenalty=10000\exhyphenpenalty=10000 \textbf{Base (famílias)}} & \parbox[b]{2.20cm}{\raggedleft\hyphenpenalty=10000\exhyphenpenalty=10000 \textbf{\% das famílias}} \\ \midrule \endhead
\midrule \multicolumn{7}{r@{}}{\footnotesize\itshape (continua)}\\ \endfoot
\bottomrule \endlastfoot
Inadequação habitacional (FJP) & 11.589 & 195.910 & 5,9 & 10.266 & 174.531 & 5,9 \\
Inadequação de infraestrutura & 9.841 & 195.622 & 5,0 & 8.696 & 174.273 & 5,0 \\
Inadequação edilícia & 2.234 & 195.910 & 1,1 & 2.005 & 174.531 & 1,1 \\
Déficit habitacional (FJP) & 86.393 & 200.108 & 43,2 & 75.681 & 178.266 & 42,5 \\
Adensamento excessivo (mais de 2 pessoas por dormitório) & 101.477 & 196.655 & 51,6 & 84.035 & 175.175 & 48,0 \\
Domicílio sem banheiro & 663 & 196.672 & 0,3 & 589 & 175.189 & 0,3 \\
Domicílio sem água canalizada & 1.213 & 196.672 & 0,6 & 1.056 & 175.189 & 0,6 \\
\end{longtable}}
\vspace{-6pt}\fonte{Elaboração IPP com dados de \citeonline{mds_cadunico}.}
\end{landscape}


A tabela resume as condições de moradia das crianças até 72 meses do CadÚnico e de suas famílias. Os problemas mais frequentes são o adensamento excessivo (51,6\% das crianças) e o déficit habitacional (43,2\%), puxado pelo peso do aluguel. A inadequação habitacional atinge 5,9\%, e a falta de banheiro (0,3\%) e de água canalizada (0,6\%) é rara. Entre as famílias, os percentuais são parecidos: 48,0\% em adensamento excessivo, 42,5\% em déficit e 5,9\% em inadequação.

\vertabelas{Tabela~\ref{tab:cadunico-deficit-componentes-0-a-5-2026}; ver Apêndice~\ref{ap:eixo-7}}

\section{Síntese do eixo}

\begin{sintese}{Moradia}O principal problema habitacional das crianças do CadÚnico não é a precariedade física da casa, que é minoritária, mas o custo do aluguel e a falta de espaço: o adensamento atinge metade delas, e o ônus excessivo com aluguel, mais de quatro em cada dez. Isso aponta para políticas de acesso à moradia e de apoio ao aluguel, além de melhorias habitacionais concentradas nos territórios com mais inadequação, como Santa Cruz e Guaratiba. Os dados das ações de melhoria habitacional do programa Territórios Sociais ainda não foram incorporados.\end{sintese}
